\section{Kramers degeneracy}
\label{sec:kramers_degeneracy}

An anti-unitary symmetry of square $-\Id$ forces every eigenvalue of a
commuting Hermitian operator to be at least two-fold degenerate.  This section
formalizes Kramers' theorem together with its antisymmetric-intertwiner
relaxation, whose printed form requires a correction
\cite[Chapter~3]{Wolf2012Quantum}.

\begin{lemma}[Independent eigenvectors give degeneracy]
    \label{lem:eigenspace_dim_of_independent_pair}
    \lean{Matrix.two_le_finrank_eigenspace_of_linearIndependent_pair}
    \leanok
    If two linearly independent vectors are eigenvectors of $M\in\MN{D}$ for the
    same eigenvalue $\lambda$, then the eigenspace of $\lambda$ has dimension at
    least two.
\end{lemma}

\begin{proof}\leanok
    The two vectors span a two-dimensional subspace of the eigenspace.
\end{proof}

\begin{theorem}[Kramers' theorem]
    \label{thm:kramers_degeneracy}
    \lean{Matrix.IsHermitian.two_le_finrank_eigenspace_of_antiunitary}
    \leanok
    Let $H\in\MN{D}$ be Hermitian and let $T$ be an anti-unitary operator with
    $[H,T]=0$ and $T^2=-\Id$.  Then every eigenvalue of $H$ is at least
    two-fold degenerate
    \cite[Chapter~3, Theorem~3.1]{Wolf2012Quantum}.
\end{theorem}

\begin{proof}\leanok
    \uses{lem:eigenspace_dim_of_independent_pair}
    Let $H\ket{\psi}=\lambda\ket{\psi}$ with $\ket{\psi}\neq 0$.  Hermiticity
    gives $\lambda\in\R$, while commutation and anti-linearity give
    \begin{align}
        H T\ket{\psi}
         & =T H\ket{\psi}
        =T\bigl(\lambda\ket{\psi}\bigr)
        =\overline\lambda T\ket{\psi}
        =\lambda T\ket{\psi}.
    \end{align}
    Antiunitarity makes $T\ket{\psi}$ nonzero.  These two eigenvectors are
    linearly independent: if $T\ket{\psi}=a\ket{\psi}$, then
    \begin{align}
        -\ket{\psi}=T^2\ket{\psi}
        =T\bigl(a\ket{\psi}\bigr)
        =\overline a a\ket{\psi},
    \end{align}
    which is impossible because $\overline a a=|a|^2\geq0$.  Hence the
    eigenspace of $\lambda$ has dimension at least two by
    Lemma~\ref{lem:eigenspace_dim_of_independent_pair}.
\end{proof}

Writing $T=\Gamma V$ with $\Gamma$ complex conjugation and $V$ unitary, the
hypotheses of Theorem~\ref{thm:kramers_degeneracy} become
$HV^\dagger=V^\dagger H^T$ and $V^T=-V$.  Wolf then relaxes the unitary to a
general antisymmetric intertwiner.  The two calculations behind the theorem
survive the relaxation verbatim; only the nonvanishing of the second
eigenvector does not (Remark~\ref{rem:kramers_ii_counterexample}).

\begin{lemma}[The partner eigenvector]
    \label{lem:kramers_partner}
    \lean{Matrix.IsHermitian.transpose_mulVec_star_of_mulVec_eq_smul,
        Matrix.IsHermitian.mulVec_star_intertwiner_of_mul_eq_mul_transpose,
        Matrix.dotProduct_star_mulVec_star_eq_zero_of_transpose_eq_neg,
        Matrix.dotProduct_mulVec_self_eq_zero_of_transpose_eq_neg}
    \leanok
    Let $H\in\MN{D}$ be Hermitian and let $A\in\MN{D}$ satisfy $HA=AH^{T}$ and
    $A^{T}=-A$.  If $H\ket{\psi}=\lambda\ket{\psi}$, then the partner vector
    $A\ket{\overline\psi}$ is again a $\lambda$-eigenvector of $H$ and is
    orthogonal to $\ket{\psi}$:
    \begin{align}
        H\,A\ket{\overline\psi}       & =\lambda\,A\ket{\overline\psi}, \\
        \braket{\psi}{A\overline\psi} & =0.
    \end{align}
\end{lemma}

\begin{proof}\leanok
    Hermiticity of $H$ makes $\lambda$ real, so conjugating the eigenvalue
    equation gives $H^{T}\ket{\overline\psi}=\lambda\ket{\overline\psi}$, and
    the intertwining relation yields
    $H\,A\ket{\overline\psi}=AH^{T}\ket{\overline\psi}
        =\lambda\,A\ket{\overline\psi}$.  For the orthogonality, antisymmetry makes
    the quadratic form of $A$ its own negative:
    $\braket{\psi}{A\overline\psi}=-\braket{\psi}{A\overline\psi}$.
\end{proof}

\begin{remark}
    \label{rem:kramers_ii_counterexample}
    Wolf prints this relaxation with the hypothesis $A\neq0$ alone and
    concludes that every eigenvalue of $H$ is at least two-fold degenerate
    \cite[Chapter~3, Theorem~3.2]{Wolf2012Quantum}.  As printed this is
    false.  On $\C^3$ with basis indexed by $0,1,2$, take
    \begin{align}
        H=\operatorname{diag}(0,1,1),\qquad
        A=
        \begin{pmatrix}
            0 & 0  & 0 \\
            0 & 0  & 1 \\
            0 & -1 & 0
        \end{pmatrix}.
    \end{align}
    Then $H$ is Hermitian, $A^{T}=-A\neq0$, and $HA=AH^{T}=A$, but the
    eigenvalue $0$ is simple: its eigenvector $\ket{e_0}$ has vanishing partner
    $A\ket{\overline{e_0}}=Ae_0=0$, so Lemma~\ref{lem:kramers_partner} supplies
    no second eigenvector.  The correct conditional statement is
    Theorem~\ref{thm:kramers_ii_nonvanishing}.
\end{remark}

\begin{theorem}[Kramers' theorem for an antisymmetric intertwiner]
    \label{thm:kramers_ii_nonvanishing}
    \lean{Matrix.IsHermitian.two_le_finrank_eigenspace_of_intertwiner_mulVec_star_ne_zero}
    \leanok
    Let $H\in\MN{D}$ be Hermitian and let $A\in\MN{D}$ satisfy $HA=AH^{T}$ and
    $A^{T}=-A$.  If $\ket{\psi}\neq0$ satisfies
    $H\ket{\psi}=\lambda\ket{\psi}$ and the partner does not vanish,
    $A\ket{\overline\psi}\neq0$, then the eigenspace of $\lambda$ has dimension
    at least two.
\end{theorem}

\begin{proof}\leanok
    \uses{lem:kramers_partner, lem:eigenspace_dim_of_independent_pair}
    By Lemma~\ref{lem:kramers_partner}, $\ket{\psi}$ and
    $A\ket{\overline\psi}$ are nonzero orthogonal eigenvectors for the same
    eigenvalue $\lambda$, hence linearly independent, and
    Lemma~\ref{lem:eigenspace_dim_of_independent_pair} gives the dimension
    bound.
\end{proof}

\begin{corollary}[Invertible intertwiner]
    \label{cor:kramers_ii_invertible}
    \lean{Matrix.IsHermitian.two_le_finrank_eigenspace_of_antisymmetric_isUnit}
    \leanok
    Let $H\in\MN{D}$ be Hermitian and let $A\in\MN{D}$ be invertible with
    $HA=AH^{T}$ and $A^{T}=-A$.  Then every eigenvalue of $H$ is at least
    two-fold degenerate.
\end{corollary}

\begin{proof}\leanok
    \uses{thm:kramers_ii_nonvanishing}
    If $H\ket{\psi}=\lambda\ket{\psi}$ with $\ket{\psi}\neq0$, invertibility
    of $A$ gives $A\ket{\overline\psi}\neq0$, so
    Theorem~\ref{thm:kramers_ii_nonvanishing} applies.  An antisymmetric
    invertible matrix exists only in even dimensions, since
    $\det A=\det A^{T}=(-1)^{D}\det A$.
\end{proof}

\begin{corollary}[Kramers' theorem, matrix form]
    \label{cor:kramers_matrix_form}
    \lean{Matrix.IsHermitian.two_le_finrank_eigenspace_of_antisymmetric_unitary,
        Matrix.transpose_conjTranspose_eq_neg_of_transpose_eq_neg}
    \leanok
    Let $H,V\in\MN{D}$ with $H$ Hermitian, $V$ unitary,
    $HV^{\dagger}=V^{\dagger}H^{T}$, and $V^{T}=-V$.  Then every eigenvalue of
    $H$ is at least two-fold degenerate.
\end{corollary}

\begin{proof}\leanok
    \uses{cor:kramers_ii_invertible}
    This is Corollary~\ref{cor:kramers_ii_invertible} at $A=V^{\dagger}$: the
    adjoint of $V$ is antisymmetric because $V$ is, and invertible because $V$
    is unitary.  This is Wolf's own proof of
    Theorem~\ref{thm:kramers_degeneracy} read through the matrix reduction.
\end{proof}
